\chapter{Ручка азарта}
% полная версия: chapters/08_nora.tex

Вы выбираете, где пообедать. Пойти в проверенное кафе или попробовать новое? Если всегда ходить в проверенное, никогда не узнаешь, что за углом есть лучше. Если всё время пробовать новое, будешь часто обедать плохо. Где золотая середина? Эта задача --- искать или пользоваться найденным --- стоит перед любым обучающимся существом. В мозге за неё, по-видимому, отвечает радиостанция \nt{НОРА} --- норадреналин.

\section*{Ручка контрастности}

Норадреналин меняет то, насколько резко нейрон отвечает на сигнал. При слабом действии нейрон похож на плавный регулятор яркости: чуть больше сигнал --- чуть ярче ответ. При сильном он похож на выключатель: ниже порога --- тишина, выше --- полный ответ. Точнее говоря, опыты показывают, что норадреналин подавляет фоновую болтовню нейронов сильнее, чем их ответ на настоящий сигнал; «ручка контрастности» --- удобная модель этого эффекта.

\pencilfig{8}{\draw[pen, wash=pcGray!50] (0,0) circle (0.65);
\draw[pen] (200:1.0) arc (200:-20:1.0);
\foreach \a in {200,165,130,95,60,25,-10} \draw[pcGray, line width=0.5pt] (\a:1.0) -- (\a:1.15);
\draw[pcRed, line width=1.2pt, line cap=round] (0,0) -- (60:0.55);
\node[lbl, anchor=east] at (200:1.25) {искать}; \node[lbl, anchor=west] at (-20:1.25) {решать};}{Ручка азарта: повернёшь влево --- машина пробует новое, вправо --- решительно выбирает проверенное.}

\section*{Думать или решать}

Вспомните из третьей главы две группы, тормозящие друг друга. Когда контрастность низкая, обе работают, слабая чуть тише: сеть ещё «думает», взвешивает варианты. Когда контрастность выше определённого порога, сеть резко переходит в режим «решила»: одна группа побеждает, другая замолкает. Одна ручка переключает машину между раздумьем и решением.

\section*{Сколько искать}

Добавим шум. При низкой контрастности шум легко перевешивает небольшую разницу между вариантами, и машина часто пробует не лучший вариант --- ищет. При высокой шум ничего не решает, и машина всегда берёт то, что считает лучшим, --- пользуется.

В нашей модели мир время от времени меняется: лучший вариант становится худшим. Выигрыш зависит от ручки как перевёрнутая буква U. Слишком мало контрастности --- машина бродит вслепую. Слишком много --- она не замечает, что мир изменился, и упрямо ходит в испортившееся кафе. И чем чаще меняется мир, тем ниже лучшая настройка: в изменчивом мире выгоднее искать.

\section*{Кто крутит ручку}

Естественный признак перемен --- неожиданность: ошибки стали больше обычного. По одной гипотезе, именно о такой неожиданной перемене и сообщает норадреналин. Тут важна тонкость: постоянный, фоновый уровень норадреналина связан скорее с поиском и рассеянностью, а короткие вспышки на важное событие --- с сосредоточенностью и решением. Возможно, вспышка работает как прерывание в компьютере: «бросай текущее, мир изменился, пересобери план».

Косвенный признак этой системы виден снаружи: зрачок расширяется, когда что-то неожиданно меняется и люди начинают учиться быстрее. Правда, зрачок следит не только за норадреналином, так что это подсказка, а не прибор.

И честное признание: в нашей модели подстройка ручки по неожиданности выиграла совсем немного по сравнению с лучшей постоянной настройкой. Где именно такая подстройка окупается, ещё предстоит понять.

\fullversion{8}
